\section{Conclusion}
This paper presented PreForceVision, a physics-guided framework for estimating bolt preload from deformation of the bolt-head front face. The framework addresses the gap between accurate but hardware-dependent preload sensing and scalable but mechanically indirect vision-based bolt inspection. Instead of treating bolt inspection as an image-classification task, the proposed approach formulates preload estimation as an inverse deformation-to-force problem. Reference-loaded image pairs are used to extract a preload-sensitive deformation signature, which is then mapped to preload under controlled conditions
%This paper presents PreForceVision, a physics-guided framework for estimating bolt preload from the deformation of the bolt-head front face. The framework addresses the gap between accurate but hardware-dependent preload sensing and scalable but mechanically indirect vision-based bolt inspection. The central idea is to compare reference and loaded images of the bolt-head face, extract a local deformation signature, and map this signature to preload under controlled conditions.

Preliminary finite-element, strain-gauge sensor integrated bolt, and optical 3D results support the mechanical basis of the concept. They indicate a measurable concave deformation of the bolt-head face under preload, with a central depression in the order of \SIrange{10}{20}{\micro\meter}. The proposed framework, therefore, focuses on measurement discipline: mechanical observability, multimodal ground truth, controlled image-pair acquisition, deformation feature extraction, physics-guided regression, and uncertainty estimation.

%Future work will implement the full experimental plan, quantify the signal-to-noise ratio across different surface preparations and camera systems, compare DIC-based baselines with learning-based models, and determine whether the approach can be extended from controlled laboratory imaging to low-cost, field-compatible acquisition.
Future work will implement the staged experimental plan, quantify the signal-to-noise ratio across surface preparations and camera systems, and compare interpretable DIC-based baselines with learning-based regression models. Particular focus on uncertainty calibration, specimen-wise validation, and robustness under controlled variations in pose, lighting, surface condition, and camera hardware. After these validation steps, the feasibility of low-cost or field-compatible acquisition can be assessed.